\section{The pressure-complete setting and mixed jet}
\label{u:sec:formal-setting-mixed-jet}

\begin{proposition}[Pressure normalization]
\label{u:prop:formal-pressure-normalization}
Let \(u,p\) be smooth on \(I\times\R^3\), let \(\nabla\cdot u=0\), and suppose
\[
 u(t)\in H^1_\sigma(\R^3),
 \qquad
 p(t)\in H^0(\R^3)
 \qquad(t\in I),
\]
and
\begin{equation}
 u_t+(u\cdot\nabla)u+\nabla p=\nu\Delta u.
 \label{u:eq:formal-ns-system}
\end{equation}
Then
\begin{align}
 -\Delta p
 &=\partial_i\partial_j(u_i u_j),
 \label{u:eq:formal-pressure-poisson}\\
 p
 &=R_iR_j(u_i u_j),
 \label{u:eq:formal-pressure-riesz}\\
 (u\cdot\nabla)u+\nabla p
 &=\Pp\nabla\cdot(u\otimes u).
 \label{u:eq:formal-pressure-projected}
\end{align}
The condition \(p(t)\in H^0(\R^3)\) fixes the pressure uniquely.
\end{proposition}

\begin{proof}
Incompressibility gives
\[
 (u\cdot\nabla)u_i=\partial_j(u_j u_i).
\]
Taking the divergence of \eqref{u:eq:formal-ns-system} cancels
\(\nabla\cdot u_t\) and \(\nabla\cdot\Delta u\), which yields
\eqref{u:eq:formal-pressure-poisson}.  Lemma~\ref{u:lem:riesz-leray}, applied
with \(m=0\), gives \(q:=R_iR_j(u_i u_j)\in H^0\), this Poisson identity, and
\[
 (u\cdot\nabla)u+\nabla q=\Pp\nabla\cdot(u\otimes u).
\]
If \(p_1,p_2\in H^0\) satisfy \eqref{u:eq:formal-pressure-poisson}, then
\(h=p_1-p_2\) satisfies
\(|\xi|^2\widehat h=0\).  Since \(\widehat h\in L^2\) and is supported at
\(\{0\}\), it vanishes.  Hence \(p=q\), which proves
\eqref{u:eq:formal-pressure-riesz} and \eqref{u:eq:formal-pressure-projected}.
\end{proof}

\begin{samepage}
\begin{definition}[Maximal pressure-complete classical history]
\label{u:def:formal-maximal-history}
Fix \(\nu>0\) and \(u_0\in\HsSigma(\R^3)\).  A pressure-complete classical
history generated by \((u_0,\nu)\) is a triple
\[
 (T_*,u,p),
 \qquad
 0<T_*\le\infty,
\]
such that, for every \(0<T<T_*\) and \(n,m\in\Nzero\),
\begin{align*}
 \partial_t^nu&\in C([0,T];H^m_\sigma(\R^3)),\\
 \partial_t^np&\in C([0,T];H^m(\R^3)),
\end{align*}
and
\begin{align}
 u_t+(u\cdot\nabla)u+\nabla p&=\nu\Delta u,
 \label{u:eq:formal-history-momentum}\\
 \nabla\cdot u&=0,
 \label{u:eq:formal-history-divergence}\\
 u(0)&=u_0,
 \label{u:eq:formal-history-datum}\\
 p(t)&=R_iR_j(u_i(t)u_j(t)).
 \label{u:eq:formal-history-pressure}
\end{align}
The history is maximal if no pressure-complete classical history
\((\widetilde T,\widetilde u,\widetilde p)\) generated by \((u_0,\nu)\)
satisfies
\[
 \widetilde T>T_*,
 \qquad
 (\widetilde u,\widetilde p)|_{[0,T_*)}=(u,p).
\]
\end{definition}
\end{samepage}

\begin{definition}[Temporal and mixed velocity--pressure jets]
\label{u:def:formal-mixed-jet}
Let \((T_*,u,p)\) be a maximal pressure-complete classical history.  For
\(n\in\Nzero\) and \(\alpha\in\Nzero^3\), define
\[
 V_n:=\partial_t^nu,
 \qquad
 Q_n:=\partial_t^np,
 \qquad
 J_{n,\alpha}:=\partial_t^n\partial_x^\alpha u,
 \qquad
 \Pi_{n,\alpha}:=\partial_t^n\partial_x^\alpha p.
\]
For \(\beta\le\alpha\), define
\[
 \binom{\alpha}{\beta}
 :=\prod_{\ell=1}^3\binom{\alpha_\ell}{\beta_\ell},
 \qquad
 |\alpha|:=\alpha_1+\alpha_2+\alpha_3.
\]
\end{definition}

\begin{proposition}[Pressure-complete binomial recurrences]
\label{u:prop:formal-mixed-recurrences}
For every \(n\in\Nzero\) and \(\alpha\in\Nzero^3\),
\begin{align}
 Q_n
 &=R_iR_j
 \sum_{a=0}^{n}\binom na V_{a,i}V_{n-a,j},
 \label{u:eq:formal-temporal-pressure-row}\\
 V_{n+1}
 &=
 \nu\Delta V_n
 -\sum_{a=0}^{n}\binom na(V_a\cdot\nabla)V_{n-a}
 -\nabla Q_n,
 \label{u:eq:formal-temporal-velocity-row}\\
 \Pi_{n,\alpha}
 &=R_iR_j
 \sum_{a=0}^{n}\sum_{\beta\le\alpha}
 \binom na\binom{\alpha}{\beta}
 J_{a,\beta,i}J_{n-a,\alpha-\beta,j},
 \label{u:eq:formal-mixed-pressure-row}\\
 J_{n+1,\alpha}
 &+
 \sum_{a=0}^{n}\sum_{\beta\le\alpha}
 \binom na\binom{\alpha}{\beta}
 (J_{a,\beta}\cdot\nabla)
 J_{n-a,\alpha-\beta}
 +\nabla\Pi_{n,\alpha}
 =\nu\Delta J_{n,\alpha},
 \label{u:eq:formal-mixed-velocity-row}\\
 \nabla\cdot J_{n,\alpha}&=0.
 \label{u:eq:formal-mixed-divergence-row}
\end{align}
\end{proposition}

\begin{proof}
Apply \(\partial_t^n\) to \eqref{u:eq:formal-history-pressure} and use
Lemma~\ref{u:lem:mixed-leibniz} with \(\alpha=0\) to obtain
\eqref{u:eq:formal-temporal-pressure-row}.  Apply \(\partial_t^n\) to
\eqref{u:eq:formal-history-momentum} to obtain
\eqref{u:eq:formal-temporal-velocity-row}.  Apply
\(\partial_x^\alpha\) to both identities, use
Lemma~\ref{u:lem:mixed-leibniz}, and commute the derivatives with \(R_iR_j\)
to obtain \eqref{u:eq:formal-mixed-pressure-row} and
\eqref{u:eq:formal-mixed-velocity-row}.  Differentiation of
\eqref{u:eq:formal-history-divergence} gives
\eqref{u:eq:formal-mixed-divergence-row}.
\end{proof}

\paragraph{The unforced datum and its derivative rows.}
The initial velocity and viscosity specify all inputs to these recurrences.
Pressure is determined by the velocity products, and there is no prescribed
force jet. A smooth forced equation can still be differentiated; its rows
also contain the corresponding force derivatives, and the pressure equation
contains the divergence of the force. In the unforced kinetic-energy
calculation of Proposition~\ref{u:prop:kinetic-energy}, transport and pressure
have zero global pairing with the velocity, and no external work term
remains. At higher Sobolev orders, the full internal transport--pressure
contribution is retained in
Proposition~\ref{u:prop:completed-height-recurrence}. These are the governing
rows and energy identities used in the construction below.

\begin{definition}[Canonical pressure constant]
\label{u:def:formal-canonical-bilinear-constants}
For vector fields \(f,g\), set
\[
 \mathcal Q(f,g):=\sum_{i,j=1}^3R_iR_j(f_i g_j).
\]
For \(m\in\Nzero\), define
\[
 C_m^{\rm p}:=
 \sup_{\substack{
 f,g\in H^{m+1}(\R^3;\R^3)\setminus\{0\}}}
 \frac{
 \norm{\mathcal Q(f,g)}_{H^m}
 +\norm{\nabla\mathcal Q(f,g)}_{H^{m-1}}}
 {\norm f_{H^{m+1}}\norm g_{H^{m+1}}}.
\]
\end{definition}

\begin{lemma}
\label{u:lem:formal-canonical-bilinear-constants}
For every \(m\in\Nzero\),
\[
 C_m^{\rm p}\le2C_m^{\rm alg}.
\]
\end{lemma}
\begin{proof}
Regard \((f_i g_j)_{i,j}\) as one tensor.  The contraction
\(F\mapsto R_iR_jF_{ij}\) has Fourier-multiplier norm one from the tensor
\(H^m\) norm to scalar \(H^m\).  Lemma~\ref{u:lem:sobolev-product} and
\(\norm{\nabla h}_{H^{m-1}}\le\norm h_{H^m}\) give
\[
 \norm{\mathcal Q(f,g)}_{H^m}
 +\norm{\nabla\mathcal Q(f,g)}_{H^{m-1}}
 \le2C_m^{\rm alg}\norm f_{H^{m+1}}\norm g_{H^{m+1}}.
\]
\end{proof}

\begin{lemma}[Initial pressure-complete jet]
\label{u:lem:formal-initial-jet}
Set
\[
 V_n^0:=V_n(0),
 \qquad
 Q_n^0:=Q_n(0),
 \qquad
 J_{n,\alpha}^0:=J_{n,\alpha}(0),
 \qquad
 \Pi_{n,\alpha}^0:=\Pi_{n,\alpha}(0).
\]
Then
\begin{align}
 V_0^0&=u_0,
 \label{u:eq:formal-initial-base}\\
 Q_n^0
 &=R_iR_j
 \sum_{a=0}^{n}\binom naV_{a,i}^0V_{n-a,j}^0,
 \label{u:eq:formal-initial-pressure}\\
 V_{n+1}^0
 &=
 \nu\Delta V_n^0
 -\sum_{a=0}^{n}\binom na
 (V_a^0\cdot\nabla)V_{n-a}^0
 -\nabla Q_n^0.
 \label{u:eq:formal-initial-velocity}
\end{align}
For every \(n,m\in\Nzero\) and \(\alpha\in\Nzero^3\),
\[
 V_n^0,Q_n^0,J_{n,\alpha}^0,\Pi_{n,\alpha}^0\in H^m(\R^3).
\]
\par
\end{lemma}

\begin{proof}
Equations \eqref{u:eq:formal-initial-base}--\eqref{u:eq:formal-initial-velocity}
are \eqref{u:eq:formal-temporal-pressure-row} and
\eqref{u:eq:formal-temporal-velocity-row} at \(t=0\).  Assume
\(V_0^0,\ldots,V_n^0\in\bigcap_{m\ge0}H^m\).  Lemmas
\ref{u:lem:sobolev-product} and \ref{u:lem:riesz-leray} give
\[
 Q_n^0\in\bigcap_{m\ge0}H^m,
 \qquad
 V_{n+1}^0\in\bigcap_{m\ge0}H^m.
\]
Induction starts from \(u_0\in\HsSigma\).  Spatial differentiation and
\eqref{u:eq:formal-mixed-pressure-row} give the claims for
\(J_{n,\alpha}^0\) and \(\Pi_{n,\alpha}^0\).
\end{proof}

\begin{theorem}[Finite-row pressure estimates]
\label{u:thm:formal-finite-row-pressure}
Let \(I=(0,T)\subset(0,T_*)\), let \(n,M\in\Nzero\), and suppose
\[
 V_a\in L^2(I;H^{M+1}(\R^3))
 \qquad(0\le a\le n).
\]
Then
\begin{align}
 &\norm{Q_n}_{L^1(I;H^M)}
 +\norm{\nabla Q_n}_{L^1(I;H^{M-1})}
 \notag\\
 &\qquad\le
 C_M^{\rm p}\sum_{a=0}^{n}\binom na
 \norm{V_a}_{L^2(I;H^{M+1})}
 \norm{V_{n-a}}_{L^2(I;H^{M+1})}.
 \label{u:eq:formal-finite-pressure-L1}
\end{align}
\begin{samepage}
If
\[
 \mathcal V_{n,M}(I)^2
 :=\sum_{a=0}^{n}
 \norm{V_a}_{L^2(I;H^{M+1})}^2,
\]
then
\begin{equation}
 \norm{Q_n}_{L^1(I;H^M)}
 +\norm{\nabla Q_n}_{L^1(I;H^{M-1})}
 \le C_M^{\rm p}2^n\mathcal V_{n,M}(I)^2.
 \label{u:eq:formal-finite-pressure-compressed}
\end{equation}
\end{samepage}
For \(\alpha\in\Nzero^3\), assume
\[
 J_{a,\beta}\in L^2(I;H^{M+1}(\R^3))
 \qquad
 (0\le a\le n,\ \beta\le\alpha).
\]
Then
\begin{align}
 &\norm{\Pi_{n,\alpha}}_{L^1(I;H^M)}
 +\norm{\nabla\Pi_{n,\alpha}}_{L^1(I;H^{M-1})}
 \notag\\
 &\quad\le
 C_M^{\rm p}
 \sum_{a=0}^{n}\sum_{\beta\le\alpha}
 \binom na\binom{\alpha}{\beta}
 \norm{J_{a,\beta}}_{L^2(I;H^{M+1})}
 \norm{J_{n-a,\alpha-\beta}}_{L^2(I;H^{M+1})}.
 \label{u:eq:formal-finite-mixed-pressure-L1}
\end{align}
If
\[
 \mathcal J_{n,\alpha,M}(I)^2
 :=
 \sum_{a=0}^{n}\sum_{\beta\le\alpha}
 \norm{J_{a,\beta}}_{L^2(I;H^{M+1})}^2,
\]
then
\begin{equation}
 \norm{\Pi_{n,\alpha}}_{L^1(I;H^M)}
 +\norm{\nabla\Pi_{n,\alpha}}_{L^1(I;H^{M-1})}
 \le
 C_M^{\rm p}2^{n+|\alpha|}
 \mathcal J_{n,\alpha,M}(I)^2.
 \label{u:eq:formal-finite-mixed-pressure-compressed}
\end{equation}
\end{theorem}

\begin{proof}
Lemmas~\ref{u:lem:sobolev-product} and \ref{u:lem:riesz-leray} give
\begin{equation}
 \norm{R_iR_j(f_i g_j)}_{H^M}
 +\norm{\nabla R_iR_j(f_i g_j)}_{H^{M-1}}
 \le C_M^{\rm p}\norm f_{H^{M+1}}\norm g_{H^{M+1}}.
 \label{u:eq:formal-pressure-bilinear}
\end{equation}
Equations \eqref{u:eq:formal-temporal-pressure-row} and
\eqref{u:eq:formal-pressure-bilinear} give, for almost every \(t\in I\),
\[
 \norm{Q_n(t)}_{H^M}
 +\norm{\nabla Q_n(t)}_{H^{M-1}}
 \le
 C_M^{\rm p}\sum_{a=0}^{n}\binom na
 \norm{V_a(t)}_{H^{M+1}}
 \norm{V_{n-a}(t)}_{H^{M+1}}.
\]
Integration and Cauchy--Schwarz prove
\eqref{u:eq:formal-finite-pressure-L1}.  For nonnegative \(x_a\),
\[
 \sum_{a=0}^{n}\binom nax_ax_{n-a}
 \le
 \sum_{a=0}^{n}\binom nax_a^2
 \le
 2^n\sum_{a=0}^{n}x_a^2.
\]
This proves \eqref{u:eq:formal-finite-pressure-compressed}.

Apply \eqref{u:eq:formal-pressure-bilinear} to each summand in
\eqref{u:eq:formal-mixed-pressure-row}.  Integration and Cauchy--Schwarz give
\eqref{u:eq:formal-finite-mixed-pressure-L1}.  The involution
\[
 (a,\beta)\longmapsto(n-a,\alpha-\beta)
\]
preserves the weights
\(\binom na\binom{\alpha}{\beta}\), and
\[
 \sum_{a=0}^{n}\sum_{\beta\le\alpha}
 \binom na\binom{\alpha}{\beta}
 =2^{n+|\alpha|}.
\]
In particular, every individual weight is at most \(2^{n+|\alpha|}\).
The inequality \(2xy\le x^2+y^2\) proves
\eqref{u:eq:formal-finite-mixed-pressure-compressed}.
\end{proof}


\subsection{One compatible rectangle family and its two exhaustions}
\label{u:sec:independent-intersection}

Start with \(u_0\in\HsSigma(\R^3)\) and \(\nu>0\).
Theorem~\ref{u:thm:local-restart} constructs a pressure-normalized classical
history on a positive interval. Uniqueness identifies the local solutions
on their overlaps; their compatible union gives the maximal history of
Definition~\ref{u:def:formal-maximal-history}.

Its derivative rows are determined by the equation. The datum specifies
the initial zeroth velocity row. The pressure formula in
Proposition~\ref{u:prop:formal-mixed-recurrences} fixes its initial pressure,
and the velocity recurrence fixes its first temporal row.
Repeating the recurrences determines
every finite initial temporal row. Spatial differentiation gives the
initial mixed jet. Along the history, those rows are the actual
distributional derivatives of \(u,p\), with the full transport and pressure
terms retained.

On every strict subinterval \((0,T)\), \(T<T_*\), ordinary persistence of
the classical history identifies all of these rows and gives finite norms at
each fixed pair of indices.  That strict-horizon statement is not the producer
of the whole-terminal result.  When a finite maximal time is proposed, put
\(I=(0,T_*)\).  The datum-generated whole-terminal theorem below is the
producing theorem.  At each fixed finite temporal depth and spatial height it
takes the datum, viscosity, complete initial face, and closed
pressure-complete equation graph, and it supplies the adjacent
whole-terminal memberships at that finite pair.  Their finitely many squared
norms form the whole-terminal rectangle norm; restriction and monotone
convergence then identify its strict-cutoff supremum.  Mixed-index exhaustion,
completed-height exhaustion, terminal tracing, and restart all occur after
that output has been obtained.

\begin{theorem}[Datum-generated whole-terminal finite rectangles]
\label{u:thm:datum-finite-blocks}
Let \(\nu>0\), \(u_0\in\HsSigma(\R^3)\), and let
\((T_*,u,p)\) be the maximal pressure-complete classical history generated
by \((u_0,\nu)\).  Assume, toward contradiction, that \(T_*<\infty\), and put
\(I=(0,T_*)\).  For finite \(N,M\), the selected pressure-complete record is
\begin{equation}
 \cR_{N,M}[u_0;T_*]
 :=\left((V_n)_{n=0}^{N},(V_{n+1})_{n=0}^{N};(Q_n)_{n=0}^{N}\right),
 \label{u:eq:independent-datum-rectangle-record}
\end{equation}
and the datum-generated construction produces, for \(0\le n\le N\),
\begin{equation}
 V_n\in L^2(I;H^{M+1}),\qquad
 V_{n+1}\in L^2(I;H^{M-1}).
 \label{u:eq:independent-adjacent-finite-rows}
\end{equation}
For \(0<S\le T_*\), define its velocity norm by
\begin{equation}
 \mathfrak R_{N,M}(u;S)
 :=\sum_{n=0}^{N}\left(
 \norm{V_n}_{L^2(0,S;H^{M+1})}^2
 +\norm{V_{n+1}}_{L^2(0,S;H^{M-1})}^2
 \right).
 \label{u:eq:independent-datum-block-output}
\end{equation}
The theorem's whole-terminal output is
\begin{equation}
 \sup_{0<S<T_*}\mathfrak R_{N,M}(u;S)
 =\mathfrak R_{N,M}(u;T_*)<\infty
 \qquad(N,M<\infty).
 \label{u:eq:datum-whole-terminal-producer}
\end{equation}
Taking \(N\ge n\) and \(M=m\) gives the corresponding statement for every
finite pair \((n,m)\).
Every row is the corresponding distributional derivative of the single
maximal history.  For every finite \(n,m,\alpha\),
\[
 J_{n,\alpha}\in L^2(I;H^m),
 \qquad \Pi_{n,\alpha}\in L^1(I;H^m),
\]
and enlargement of a finite rectangle preserves every common velocity and
pressure coordinate.
\end{theorem}

\begin{proof}[Datum construction and finite-coordinate realization]
Put \(I=(0,T_*)\), and fix finite \(N,M\).  The datum first generates the
complete initial jet
\[
 A_0=u_0,
 \qquad
 P_n=R_iR_j\sum_{a=0}^n\binom na(A_a)_i(A_{n-a})_j,
\]
\[
 A_{n+1}=\nu\Delta A_n
 -\sum_{a=0}^n\binom na(A_a\cdot\nabla)A_{n-a}-\nabla P_n.
\]
At step \(n\), the right-hand side uses only the already determined rows
\(A_0,\ldots,A_n\).  Lemma~\ref{u:lem:formal-initial-jet} therefore places
every required \(A_n\) and \(P_n\) in every finite spatial Sobolev order.
Spatial differentiation supplies
\[
 J_{n,\alpha}(0)=\partial_x^\alpha A_n,
 \qquad
 \Pi_{n,\alpha}(0)=\partial_x^\alpha P_n.
\]

The selected equations retain their complete finite dependency cover.  For
an equation indexed by \((n,\alpha)\), retain the adjacent temporal coordinate
\(J_{n+1,\alpha}\), the viscous coordinates
\(J_{n,\alpha+2e_\ell}\), and the pressure coordinates
\(\Pi_{n,\alpha+e_\ell}\), for \(1\le\ell\le3\).  For every \(a\le n\),
\(\beta\le\alpha\), and velocity component \(i\), retain both transport
factors
\[
 J_{a,\beta,i},
 \qquad
 J_{n-a,\alpha-\beta+e_i}.
\]
Here \(e_i\) is the \(i\)-th spatial unit multi-index.  Set
\(\Gamma_\alpha:=\{\alpha\}\cup\{\alpha+e_\ell:1\le\ell\le3\}\).
For every \(\gamma\in\Gamma_\alpha\), every \(a\le n\), and every
\(\beta\le\gamma\), also retain the two pressure-product factors
\(J_{a,\beta,i}\) and \(J_{n-a,\gamma-\beta,j}\).  Each retained normalized
pressure coordinate then carries its full product formula
\[
 \Pi_{n,\gamma}
 =R_iR_j\sum_{a=0}^n\sum_{\beta\le\gamma}
 \binom na\binom{\gamma}{\beta}
 J_{a,\beta,i}J_{n-a,\gamma-\beta,j}
 \qquad(\gamma\in\Gamma_\alpha).
\]
Thus the retained equation is
\begin{equation}
 \begin{aligned}
 \mathscr E_{n,\alpha}(J,\Pi):={}&J_{n+1,\alpha}
 +\sum_{a=0}^{n}\sum_{\beta\le\alpha}
 \binom na\binom{\alpha}{\beta}
 (J_{a,\beta}\cdot\nabla)J_{n-a,\alpha-\beta}\\
 &+\nabla\Pi_{n,\alpha}-\nu\Delta J_{n,\alpha}=0.
 \end{aligned}
 \label{u:eq:finite-equation-graph}
\end{equation}
There are finitely many retained coordinates for the fixed selected equations.
The dependency cover is not an autonomous truncated fluid: every entry is the
corresponding derivative of the original maximal history.  Applying \(\Pp\)
reads the solenoidal evolution, and applying \(I-\Pp\) reads the normalized
pressure equation.  The initial face, every nonlinear factor, the viscous
coordinates, and the pressure coordinates are therefore present together.

The datum-generated columnwise construction now acts on this one complete
finite record.  For each \(n=0,\ldots,N\), it supplies
\[
 V_n\in L^2(I;H^{M+1}),
 \qquad
 V_{n+1}\in L^2(I;H^{M-1}).
\]
These are the whole-terminal memberships in
\eqref{u:eq:independent-adjacent-finite-rows}; they precede every cutoff norm
and every terminal trace.  The adjacent entry is the weak temporal derivative
of the preceding entry by \eqref{u:eq:finite-equation-graph}.

For this fixed finite pair, form the sum of the defining squared norms:
\begin{equation}
 C_{N,M}[u_0,\nu;I]
 :=\sum_{n=0}^{N}\left(
  \norm{V_n}_{L^2(I;H^{M+1})}^2
  +\norm{V_{n+1}}_{L^2(I;H^{M-1})}^2
 \right)<\infty.
 \label{u:eq:finite-graph-output}
\end{equation}
Each displayed membership asserts finiteness of its defining squared norm,
and the sum has only finitely many terms.
Restriction to \((0,S)\) decreases each nonnegative integral.  Increasing
\(S\) to \(T_*\) and applying monotone convergence to the finite sum gives
\begin{equation}
 \sup_{0<S<T_*}\mathfrak R_{N,M}(u;S)
 =\mathfrak R_{N,M}(u;T_*)
 =C_{N,M}[u_0,\nu;I]<\infty.
 \label{u:eq:finite-graph-uniform-cutoff-output}
\end{equation}
The constant belongs to the selected finite pair; no constant common to the
infinite tower is asserted.  This proves
\eqref{u:eq:datum-whole-terminal-producer} in the same order as the
construction: memberships, finite norm, restriction, monotone terminal
realization.

Taking \(N\ge n\) and \(M=m\) gives the two memberships in
\eqref{u:eq:independent-adjacent-finite-rows}.  Taking
\(M\ge m+|\alpha|\) and differentiating spatially gives
\(J_{n,\alpha}\in L^2(I;H^m)\).  The pressure formula and
Theorem~\ref{u:thm:formal-finite-row-pressure} give
\[
 \norm{\Pi_{n,\alpha}}_{L^1(I;H^m)}
 \le C_{n,m,\alpha}
 \sum_{a=0}^{n}
 \norm{V_a}_{L^2(I;H^{m+|\alpha|+1})}
 \norm{V_{n-a}}_{L^2(I;H^{m+|\alpha|+1})}<\infty.
\]

When \(N'\ge N\) and \(M'\ge M\), both rectangles begin with the same initial
jet.  The temporal recurrence, spatial Leibniz rule, and normalized pressure
formula give the same expressions on every common coordinate; uniqueness
identifies their realizations.  Deleting the outer coordinates therefore
recovers the smaller rectangle.  Time restriction preserves those same
coordinates, while \eqref{u:eq:finite-graph-output} and
\eqref{u:eq:finite-graph-uniform-cutoff-output} identify their
whole-terminal norms and strict-cutoff realizations.
\end{proof}

Its two cofinal readings are consequences of the same finite-row output,
rather than consequences of one another.

\begin{corollary}[Mixed-index exhaustion of the compatible rectangles]
\label{u:thm:datum-intersection}
Under the hypotheses of Theorem~\ref{u:thm:datum-finite-blocks}, compatibility
of the finite rows gives
\begin{equation}
 u\in\bigcap_{r,m\in\Nzero}H^r(I;H^m_\sigma(\R^3)),
 \qquad I=(0,T_*).
 \label{u:eq:independent-mixed-intersection}
\end{equation}
Its mixed velocity--pressure coordinates are
\begin{equation}
 J_{n,\alpha}=\partial_x^\alpha\partial_t^nu,
 \quad
 \Pi_{n,\alpha}=\partial_x^\alpha\partial_t^np,
 \qquad n\in\Nzero,\quad\alpha\in\Nzero^3.
 \label{u:eq:independent-intersection-record}
\end{equation}
These coordinates obey the pressure-complete recurrences of
Proposition~\ref{u:prop:formal-mixed-recurrences} and agree under every
temporal and spatial restriction.
\end{corollary}

\begin{proof}
Fix finite \(r,m\).  A rectangle with temporal depth at least \(r\) and spatial
height at least \(m\) contains \(V_n\in L^2(I;H^m)\) for
\(0\le n\le r\).  The recurrence identifies these rows with the weak
derivatives \(\partial_t^nu\), so
\[
 \norm u_{H^r(I;H^m)}^2
 =\sum_{n=0}^{r}\norm{V_n}_{L^2(I;H^m)}^2<\infty.
\]
Taking every finite \((r,m)\) proves
\eqref{u:eq:independent-mixed-intersection}.  The coordinate memberships and
restriction compatibility in Theorem~\ref{u:thm:datum-finite-blocks} give
\eqref{u:eq:independent-intersection-record}.
\end{proof}

The return from the spatial column to the mixed jet uses only the original
equation and the one fixed history.

\begin{proposition}[Spatial-column reconstruction of the mixed intersection]
\label{u:prop:spatial-column-reconstruction}
Let \((T_*,u,p)\) be the pressure-normalized maximal history with
\(T_*<\infty\), put \(I=(0,T_*)\), and assume
\begin{equation}
 u\in\mathscr S_0(I):=\bigcap_{m\in\Nzero}L^2(I;H^m_\sigma(\R^3)).
 \label{u:eq:complete-spatial-row}
\end{equation}
Then \(V_n,Q_n\in C([0,T_*];H^m)\) for every finite \(n,m\), the adjacent
rows satisfy \eqref{u:eq:independent-adjacent-finite-rows}, and
\eqref{u:eq:independent-mixed-intersection} follows for this same history.
\end{proposition}

\begin{proof}
Fix \(m\). Lemmas~\ref{u:lem:sobolev-product} and \ref{u:lem:riesz-leray}
give
\[
 \begin{aligned}
 \norm{\Pp\nabla\!\cdot(u\otimes u)}_{L^1(I;H^m)}
 &\le C_m\norm u_{L^2(I;H^{m+2})}^2<\infty,\\
 \norm{\Delta u}_{L^1(I;H^m)}
 &\le\sqrt{T_*}\norm u_{L^2(I;H^{m+2})}<\infty.
 \end{aligned}
\]
Substitution into the unforced equation
\[
 u_t=\nu\Delta u-\Pp\nabla\!\cdot(u\otimes u)
\]
gives \(u_t\in L^1(I;H^m)\). Also \(u\in L^1(I;H^m)\), since
\(T_*<\infty\). The Banach-valued fundamental theorem gives
\[
 u\in W^{1,1}(I;H^m)\hookrightarrow C([0,T_*];H^m)
 \qquad(m\in\Nzero).
\]
Integrating the equation from the initial datum makes the trace bound explicit:
\[
 \sup_{0\le t\le T_*}\norm{u(t)}_{H^m}
 \le\norm{u_0}_{H^m}
 +\nu\sqrt{T_*}\norm u_{L^2(I;H^{m+2})}
 +C_m\norm u_{L^2(I;H^{m+2})}^2<\infty.
\]
These continuous representatives agree under Sobolev embedding because
they extend the original velocity. Their endpoint values define one
\(u_*\in H^\infty_\sigma\).

Now assume the temporal rows \(V_0,\ldots,V_n\) are continuous on
\([0,T_*]\) in every finite Sobolev space. This has just been proved
for \(V_0=u\). The pressure and velocity recurrences are
\[
 \begin{aligned}
 Q_n&=R_iR_j\sum_{a=0}^n\binom na(V_a)_i(V_{n-a})_j,\\
 V_{n+1}&=\nu\Delta V_n
 -\sum_{a=0}^n\binom na(V_a\cdot\nabla)V_{n-a}-\nabla Q_n.
 \end{aligned}
\]
Products, spatial differentiation, and the Riesz operators are continuous
between the required finite Sobolev spaces. Choosing sufficiently high
spatial orders in the available rows makes both right-hand sides continuous
in any prescribed \(H^m\). On \(I\), they equal the actual derivatives
\(Q_n=\partial_t^np\) and \(V_{n+1}=\partial_tV_n\) by
Proposition~\ref{u:prop:formal-mixed-recurrences}. The identity
\[
 V_n(t)-V_n(s)=\int_s^tV_{n+1}(\tau)\,d\tau
 \qquad\text{in }H^m
\]
extends from interior times to \(0\le s\le t\le T_*\) by continuity.
It identifies the continuous extension with the temporal derivative,
including the one-sided derivatives at the endpoints. Induction gives
\(V_n,Q_n\in C([0,T_*];H^m)\) for every finite \(n,m\).
Finite \(T_*\) now gives the corresponding \(L^2(I;H^m)\) memberships.
In particular, selecting adjacent temporal coordinates gives
\[
 V_n\in L^2(I;H^{m+1}),\qquad
 V_{n+1}\in L^2(I;H^{m-1})
 \qquad(n,m\in\Nzero).
\]
For fixed finite \(r,m\), these actual weak derivatives give
\[
 \norm u_{H^r(I;H^m)}^2
 =\sum_{n=0}^r\norm{V_n}_{L^2(I;H^m)}^2<\infty.
\]
Thus the full-interval spatial component propagates through
the original equation to every finite mixed Sobolev membership. The argument
uses finitely many temporal orders at each step, with no bound required to
be uniform over \(r,m\). The identities in
Proposition~\ref{u:prop:formal-mixed-recurrences} give compatibility:
restricting the indices retains the corresponding derivatives of \(u,p\).
These are the finite coordinates used in
Theorem~\ref{u:thm:datum-rectangles} and the derivative-sum calculation of
Theorem~\ref{u:thm:compatible-inverse-intersection}.
\end{proof}

The adjacent memberships also give a bound valid up to the proposed terminal
time. Apply Lemma~\ref{u:lem:adjacent-trace} with
\(\partial_tV_n=V_{n+1}\). It gives
\[
 V_n\in C([0,T_*];H^m)
\]
and, using the initial value determined in
Lemma~\ref{u:lem:formal-initial-jet},
\[
 \sup_{0\le t\le T_*}\norm{V_n(t)}_{H^m}^2
 \le\norm{V_n(0)}_{H^m}^2
 +2\norm{V_{n+1}}_{L^2(I;H^{m-1})}
    \norm{V_n}_{L^2(I;H^{m+1})}<\infty.
\]
Both time norms on the right are supplied on the entire interval \(I\),
so the bound remains fixed as \(t\) approaches \(T_*\).
The trace at each spatial order is a continuous representative of the
original row. Uniqueness of these representatives under Sobolev embedding
makes their terminal values compatible. The detailed endpoint calculation
in Theorem~\ref{u:thm:ier-compatible-endpoint-traces} uses precisely these
rows and norms.

The pressure spaces in the statement retain the original normalization.
For fixed \(n,m\), the compatible finite-row family gives
\(V_a\in L^2(I;H^{m+1})\) for \(0\le a\le n\).
The pressure formula and the finite-row product estimate give
\[
 \norm{Q_n}_{L^1(I;H^m)}
 \le C_m^{\rm p}\sum_{a=0}^n\binom na
 \norm{V_a}_{L^2(I;H^{m+1})}
 \norm{V_{n-a}}_{L^2(I;H^{m+1})}<\infty.
\]
Choosing the spatial order higher by \(|\alpha|\) gives the asserted
mixed pressure coordinate.  Differentiation and restriction act on the
fixed distributions \(u,p\), so their overlapping coordinates agree.
This pressure completion is already carried by the compatible finite rows of
Theorem~\ref{u:thm:datum-finite-blocks}; it is retained in both descriptions.

Fixing temporal and spatial orders selects a finite rectangle from the mixed
jet.  Adding its squared row norms gives Theorem~\ref{u:thm:datum-rectangles},
and monotone convergence gives its value on \((0,T_*)\) in
Proposition~\ref{u:prop:datum-terminal-realization}.  Letting the rectangle
indices increase exhausts every finite mixed velocity and pressure row.  The next
section derives the completed-height exhaustion on the same finite rectangles:
all finite \(R_k\) give the complete spatial Sobolev column, and the recurrence
above reconstructs every temporal and pressure row.  The two routes meet at
the same direct mixed Sobolev membership.

The endpoint argument uses adjacent temporal coordinates of this family.
Theorem~\ref{u:thm:ier-compatible-endpoint-traces} applies the Hilbert-scale
trace estimate to obtain their compatible limits at \(T_*\).
Theorem~\ref{u:thm:ier-pressure-complete-endpoint} passes the pressure formula
and velocity recurrence to those limits. The resulting smooth endpoint
and its jet are the data used in
Theorem~\ref{u:thm:ier-same-history-restart}; local existence and overlap
uniqueness continue the original history past the proposed finite maximal
time. The small-critical-data calculation in
Theorem~\ref{u:thm:app-small-critical-global} gives a separate quantitative
result under its stated smallness condition.
